\documentclass[11pt,letterpaper]{article}
\usepackage[margin=1.04in,top=0.83in,bottom=0.85in]{geometry}
\usepackage{amsmath,amssymb,booktabs,xcolor}
\usepackage[colorlinks=true,urlcolor=accent,linkcolor=accent,hypertexnames=false]{hyperref}
\definecolor{accent}{RGB}{32,79,91}
\hypersetup{pdftitle={A certified Fourier bound for the least quadratic nonresidue}}
\setlength{\parindent}{0pt}
\setlength{\parskip}{5pt plus 0.5pt minus 0.5pt}
\raggedbottom
\setlength{\abovedisplayskip}{6pt plus 1pt}
\setlength{\belowdisplayskip}{6pt plus 1pt}
\setlength{\abovedisplayshortskip}{4pt plus 1pt}
\setlength{\belowdisplayshortskip}{4pt plus 1pt}
\newcommand{\head}[1]{\par\vspace{5pt}{\large\sffamily\bfseries\color{accent}#1}\par\nobreak}
\newcommand{\R}{\mathbb R}
\newcommand{\Fhat}{\widehat F}
\newcommand{\norm}[1]{\lVert #1\rVert}
\newcommand{\running}[1]{{\small\sffamily\color{accent}#1}\par\vspace{1pt}\hrule\vspace{9pt}}
\begin{document}
{\small\sffamily\color{accent}Working draft \(\mathbin{\bullet}\) November 20, 2026}\par
\vspace{3pt}
{\LARGE\bfseries A certified Fourier bound for the\\[2pt]
least quadratic nonresidue\par}
\vspace{4pt}
An explicit test function with 36 rational coefficients gives
\[
                         \boxed{\,\mathcal C(1)^{-2}<0.760777412\,}
\]
for the Fourier extremal constant of Carneiro, Milinovich, Quesada-Herrera, and Ramos~[1].
Their published upper bound is \(0.7615\). The numerical inequality above is
unconditional. Under GRH for Dirichlet \(L\)-functions it also bounds the
asymptotic coefficients for the least quadratic nonresidue and the least prime
quadratic residue modulo a prime.
\head{1\quad The extremal quotient}
Use \(\Fhat(t)=\int_{\R}F(x)e^{-2\pi ixt}\,dx\). For nonzero \(F\in L^1(\R)\)
with real-valued Fourier transform, write
\begin{equation}
 J(F)=\frac{2\pi}{\norm F_1}\left(\int_{-\infty}^0 e^{\pi t}\Fhat(t)\,dt
               -\int_0^\infty e^{\pi t}|\Fhat(t)|\,dt\right).
 \tag{1}
\end{equation}
The supremum of \(J(F)\) over this class is \(\mathcal C(1)\) by [1, (1.1)].
Any witness with \(J(F)>0\) gives \(\mathcal C(1)^{-2}\le J(F)^{-2}\).

\head{2\quad A single-shift test function}
Set \(a=3/10\) and \(c=4/5\). Let \(R(z)=\sum_{j=0}^{35}r_jz^j\), with the exact
coefficients on page~3. Define \(d_k\) and \(P\) by
\[
 R(2w-1)=\sum_{k=1}^{36}d_k w^{k-1},
 \qquad P(u)=\sum_{k=1}^{36}d_k\,\frac{u^k}{k!}.
\]
Take \(\Fhat(t)=\mathbf 1_{\{u\ge0\}}e^{-u}P(u)\), where \(u=(c-\pi t)/a\).
This transform is real, continuous at the cutoff because \(P(0)=0\), and
integrable. Its two weighted integrals in (1) converge absolutely. Substitution
into the inverse Fourier integral gives
\[
 F(x)=\frac a\pi e^{2icx}\sum_{k=1}^{36}
      \frac{d_k}{k!}\int_0^\infty e^{-(1+2iax)u}u^k\,du
     =\frac a\pi e^{2icx}\sum_{k=1}^{36}\frac{d_k}{(1+2iax)^{k+1}}.
\]
Thus \(F(x)=O((1+|x|)^{-2})\). Fourier inversion recovers the prescribed
continuous transform, and \(R\ne0\) ensures \(F\ne0\). This proves admissibility.

To compute the norm, use evenness of \(|F|\) and put
\(\theta=\arctan(2ax)\) for \(x\ge0\).
Writing \(w=(1+i\tan\theta)^{-1}=\cos\theta\,e^{-i\theta}\) gives
\(2w-1=e^{-2i\theta}\). The factor \(w^2\) in the sum for \(F\) has modulus
\(\cos^2\theta\), which cancels the \(\sec^2\theta\) in \(2a\,dx\). Hence
\begin{equation}
 \norm F_1=\frac{D}{\pi},\qquad
 D=\int_0^{\pi/2}|R(e^{-2i\theta})|\,d\theta. \tag{2}
\end{equation}
\newpage
\running{A certified Fourier bound}
\head{3\quad Evaluation of the signed numerator}
Put \(h=1+a=13/10\) and \(m=c/a=8/3\). In (1), \(t<0\) corresponds
to \(u>m\), while \(0<t<c/\pi\) corresponds to \(0<u<m\). Using (2) and
\(e^{\pi t}=e^{c-au}\) yields \(J(F)=N/D\), where
\begin{equation}
 N=2\pi a e^c\left(\int_m^\infty e^{-hu}P(u)\,du-
                    \int_0^m e^{-hu}|P(u)|\,du\right). \tag{3}
\end{equation}
For \(v\ge0\) the signed tail has the exact expression
\[
 T(v)=\int_v^\infty e^{-hu}P(u)\,du
     =e^{-hv}\sum_{j=0}^{36}\frac{v^j}{j!}
                    \sum_{k=j}^{36}\frac{d_k}{h^{k-j+1}},
 \qquad d_0=0,
\]
by repeated integration by parts. Exact rational real-root isolation of \(P\)
on \([0,m]\) gives the simple root at zero and two simple interior roots.
Their locations, rounded out for display, are
\[
 0.31184795095143<v_1<0.31184795095144,\qquad
 1.05793812431641<v_2<1.05793812431642.
\]
The verifier isolates each to an interval of width at most \(10^{-30}\)
and includes both endpoints in the subsequent ball arithmetic. Put \(v_0=0\),
\(v_3=m\). On every open interval \((v_j,v_{j+1})\), \(P\) has constant sign.
Consequently the second integral in (3) is exactly
\(\sum_{j=0}^{2}|T(v_j)-T(v_{j+1})|\).

\head{4\quad A one-sided integration error}
Let \(f(\theta)=R(e^{-2i\theta})\), \(L=\pi/2\), \(n=32768\), \(b=L/n\).
Differentiating the trigonometric polynomial twice gives
\(\sup|f''|\le K=4\sum_{j=0}^{35} j^2|r_j|=126726777/50000\).
On a panel \([s,s+b]\), let \(\ell\) interpolate \(f\) linearly at the
endpoints. At \(s+t\) its error satisfies
\[
 |f(s+t)-\ell(s+t)|\le K\,t(b-t)/2\qquad (0\le t\le b).
\]
Indeed the interpolation error is an integral of \(f''\) against the signed
Dirichlet Green kernel. Its absolute kernel is
\(\min(t,y)(b-\max(t,y))/b\) at \(0\le y\le b\), whose integral in \(y\) is
\(t(b-t)/2\). Convexity also gives
\(\int_s^{s+b}|\ell|\le b(|f(s)|+|f(s+b)|)/2\). Summing both bounds gives
\begin{equation}
 D\le B_n+\frac{KLb^2}{12},\qquad
 B_n=b\left(\frac{|f(0)|+|f(L)|}{2}+\sum_{j=1}^{n-1}|f(jb)|\right).
 \tag{4}
\end{equation}
This estimate remains valid at zeros of \(f\), since it bounds the interpolation
error of the complex polynomial before taking its modulus.

\newpage
\running{A certified Fourier bound}
\head{5\quad Exact coefficients and numerical certificate}
The entries \(q_j\) in the table are integers. Set \(r_j=q_j/10^6\). Thus
the index \(j\) gives the power of \(z\) in \(R(z)\).

\begin{center}
\setlength{\tabcolsep}{10pt}\renewcommand{\arraystretch}{1.13}
\begin{tabular}{@{}rr@{\qquad}rr@{\qquad}rr@{\qquad}rr@{}}
\toprule
\(j\)&\(q_j\)&\(j\)&\(q_j\)&\(j\)&\(q_j\)&\(j\)&\(q_j\)\\
\midrule
0 & 1000000 &  9 & 510409 & 18 & 6332    & 27 & \(-1653\)\\
1 & 4815715 & 10 & 258500 & 19 & 5339    & 28 & 1541\\
2 & 8710552 & 11 & 235574 & 20 & \(-4625\) & 29 & \(-1050\)\\
3 & 7617483 & 12 & 72629  & 21 & 6832    & 30 & 1019\\
4 & 4487466 & 13 & 71024  & 22 & \(-4287\) & 31 & \(-464\)\\
5 & 3434816 & 14 & 46027  & 23 & 2822    & 32 & 1320\\
6 & 2345152 & 15 & 3286   & 24 & \(-479\)  & 33 & 215\\
7 & 1189056 & 16 & 28100  & 25 & \(-732\)  & 34 & 2070\\
8 & 938349  & 17 & \(-2855\) & 26 & 1601 & 35 & 5064\\
\bottomrule
\end{tabular}
\end{center}
The script \texttt{verify\_single\_shift\_36\_arb.py} forms the rational polynomials
and isolates the roots with SymPy. It encloses all remaining real and complex
operations with 200-bit Arb balls through \texttt{python-flint}. The following
inequalities have been rounded outward.
\[
\begin{aligned}
 N &>16.68808877572884,\\
 B_n &<14.55577666588763,\\
 KLb^2/12 &<0.000000762389490,\\
 D &<14.55577742827712.
\end{aligned}
\]
Combining (1)--(4) gives
\[
 \mathcal C(1)^{-2}\le (D/N)^2
 \le \left((B_n+KLb^2/12)/N\right)^2<0.760777411894<0.760777412.
\]
A second verifier uses Cauchy--Schwarz on 2048 panels, integrating
\(|f|^2\) on each panel as a finite trigonometric sum. Applied to the same
rational polynomial, this second denominator estimate gives \(\mathcal C(1)^{-2}<0.760778\).

\head{Arithmetic consequence}
Assume GRH for Dirichlet \(L\)-functions. For each odd prime \(p\), let \(n_p\)
be its least positive quadratic nonresidue and \(r_p\) its least prime quadratic
residue. By Theorems 1 and 2 of [1], both
\(\limsup_{p\to\infty}n_p/(\log p)^2\) and
\(\limsup_{p\to\infty}r_p/(\log p)^2\) are at most \(\mathcal C(1)^{-2}\).
The displayed certificate therefore bounds each limsup strictly by \(0.760777412\).

\head{Reference}
{\small [1] E.~Carneiro, M.\,B.~Milinovich, E.~Quesada-Herrera, and A.\,P.~Ramos,
\emph{Fourier optimization, the least quadratic non-residue, and the least prime
in an arithmetic progression}, Mathematics of Computation.
\href{https://doi.org/10.1090/mcom/4154}{doi:10.1090/mcom/4154}\enspace
\href{https://arxiv.org/abs/2404.08380}{arXiv:2404.08380}.}
\end{document}
